% !TEX root = ../../../MathModel.tex
% 负责人：罗财能。
\subsubsection{具体分析}
\label{sec:q4:analysis}

问题四研究固定救援任务在分组独立执行后需要多少资源。
问题三的货箱组批、运输机型、访问顺序、准备与返航时刻、中继悬停位置及服务窗、通信保障关联全部保持不变，不能通过重新排序或错峰出发降低资源需求。
在此条件下，货箱交付时刻、运输与中继架次数以及总能耗均不因分组改变；发生变化的是同类资源可以复用的任务范围。
原本先后执行、可共用一架无人机或一组电池的任务，分属不同任务组后必须分别配置资源。

本文将“任务安排不变”解释为任务内容与时序不变，同型实体无人机和能源资源编号可以重新分配。
这使本问能够核算独立执行的最少资源，而非机械统计第三题出现过的编号。
重新分配仅发生在各组内部，不允许同一实体在执行期间跨组服务。
题目直接规定，同一运输架次涉及的服务区必须同组。
对于中继，本文采用不拆分、不复制原有中继架次的口径；结合通信保障关系不变及资源不得跨组调配，进一步要求同一中继架次保障的服务区同组。
后一要求是本文解释下的推论，不是题目额外明列的约束。

因此，应先合并不能拆开的服务区，再比较这些单元的两组和三组划分。
各组资源需求由固定占用区间的最大重叠数量确定，既包括执行任务，也包括充电或中继机周转造成的暂不可用时间。
题目未要求任务组在地理上连通，本文不另加空间连通约束。
评价时分别考察库存缺口、资源配置总量和工作量均衡，避免将资源较少直接等同于分工合理。
